\chapter{Stand der Technik}

Während das Grundlagenkapitel etabliertes Wissen vermittelt, beschreibt der Stand der Technik den \textit{aktuellen} Wissensstand. Hier wird gezeigt, welche Lösungsansätze bereits existieren, wo Forschungslücken bestehen und wie sich die eigene Arbeit davon abgrenzt.

Typische Quellen sind Journalartikel, Konferenzbeiträge und technische Berichte — also aktuelle, begutachtete Veröffentlichungen.

\section{Literaturrecherche}

Fachartikel aus begutachteten Zeitschriften (Journals) gelten als die zuverlässigsten wissenschaftlichen Quellen. Sie durchlaufen ein Peer-Review-Verfahren und sind in der Regel über eine DOI dauerhaft abrufbar.

Neben Journalartikeln spielen auch Berichte renommierter Institutionen eine wichtige Rolle, insbesondere bei der Einordnung von Zukunftstechnologien. Technische Berichte werden mit \texttt{@techreport} in der \texttt{.bib}-Datei erfasst.

\section{Verwandte Arbeiten}

Konferenzbeiträge bieten häufig den aktuellsten Einblick in laufende Forschung, da sie schneller veröffentlicht werden als Journalartikel.

\section{Abgrenzung und Forschungslücke}

Am Ende dieses Kapitels sollte klar herausgearbeitet werden, welche Lücke die eigene Arbeit schließt. Eine bewährte Struktur ist: (1) Was wurde bereits erforscht? (2) Was fehlt noch? (3) Wie adressiert die eigene Arbeit genau diese Lücke? Dieser Abschnitt bildet die Brücke zwischen der Literatur und der eigenen Methodik.